\documentclass{article}

\usepackage[utf8]{inputenc}
\usepackage[T1]{fontenc}
\usepackage{amsmath,amssymb,amsthm}
\usepackage{mathtools}
\usepackage{hyperref}
\usepackage{booktabs}

\makeatletter
\newcommand{\citep}[2][]{%
  \begingroup
  \def\@tmpkey{#2}%
  \ifx\@tmpkey\@empty\else
    (\ifx\@empty#1\@empty\else #1, \fi
     \@nameuse{bibkey@\@tmpkey})%
  \fi
  \endgroup
}
\expandafter\def\csname bibkey@lasser2026exo\endcsname{Exogenous~Verification}
\expandafter\def\csname bibkey@lasser2026phase\endcsname{Phase~Redundancy}
\expandafter\def\csname bibkey@lasser2026micro\endcsname{Microfoundation}
\makeatother

\newtheorem{definition}{Definition}
\newtheorem{lemma}{Lemma}
\newtheorem{theorem}{Theorem}
\newtheorem{remark}{Remark}
\newtheorem{assumption}{Assumption}

\newcommand{\Aadv}{A_{\mathrm{adv}}}
\newcommand{\deltaadv}{\delta_{\mathrm{adv}}}
\newcommand{\Pproxy}{P}
\newcommand{\Ttruth}{T}
\newcommand{\epsgap}{\varepsilon_{\mathrm{gap}}}
\newcommand{\epsnonres}{\varepsilon_{\mathrm{gap}}^{\mathrm{nonres}}}
\newcommand{\epsfloor}{\varepsilon_{\mathrm{floor}}^{\mathrm{res}}}

\title{Lemma 6 ($h_{\mathrm{detect}}$ intensivity): Full Proof Draft}
\author{Paper 10 working draft for codex review}
\date{}

\begin{document}

\maketitle

\section{Goal of this draft}

Theorem~\ref{thm:deployment_safety}'s Layer~2 (detection with lead
time for correction) asserts a bound during the detection window:
\[
  |g(\Ttruth) - g(\Pproxy)| \;\leq\; \mathrm{Lip}(g) \cdot
  h_{\mathrm{detect}}(\theta),
\]
with $h_{\mathrm{detect}}$ ``also intensive in $|P|$.''  V1 of the
theorem proof (Step~6 of Layer~2) referenced
\citep[Microfoundation]{lasser2026micro}'s ``$T$ failure modes''
remark for the per-step gap-growth rate, but did not formalize
either the rate's intensivity or the integrated-window bound
$h_{\mathrm{detect}}$.

This draft promotes the assertion to a formal Lemma~6, with an
explicit (C11) bounded-gap-growth-rate condition added to
Theorem~\ref{thm:deployment_safety}.

\section{Setup}

During the detection window $[t_{\mathrm{violation}},
t_{\mathrm{violation}} + T_{\mathrm{detect}}]$, an adversarial
strategy in $\Aadv \cup \Aadv^{\mathrm{env}}$ produces channel
deviations bounded by Lemma~5b's KL floor $\deltaadv$ (else SPRT
would have already detected).  The proxy-truth gap may grow during
this window before SPRT detects and operators correct.

The integrated gap during the window:
\[
  h_{\mathrm{detect}}(\theta) \;:=\;
  \int_{t_{\mathrm{violation}}}^{t_{\mathrm{violation}} +
  T_{\mathrm{detect}}}
  \frac{d \epsgap}{dt} \, dt
  \;\leq\;
  T_{\mathrm{detect}} \cdot \sup_{t \in
  [t_{\mathrm{violation}}, t_{\mathrm{violation}} +
  T_{\mathrm{detect}}]}
  \frac{d \epsgap}{dt}.
\]
For Layer~2's intensivity claim, both factors must be intensive in
$|P|$ over deployment class $\mathcal{D}$.

By Lemma~4 (SPRT lead-time tail bound), $T_{\mathrm{detect}}$ has
exponentially-decaying tail with constants intensive in $|P|$;
specifically, $\mathbb{E}[T_{\mathrm{detect}}] \leq A/\deltaadv$
and $\Pr[T_{\mathrm{detect}} > T_{\mathrm{cascade}}] \leq \beta'$
for a $|P|$-independent bound.  We treat $T_{\mathrm{detect}}$ as
intensive over $\mathcal{D}$ in the sense that operators set $A$
and $\deltaadv$ as deployment-class constants.

The remaining piece is the per-step (or per-unit-time) gap-growth
rate $d\epsgap/dt$.  This is what (C11) bounds.

\section{Bounded-gap-growth-rate condition}

\begin{assumption}[(C11): Bounded per-unit-time gap-growth rate during detection window]
\label{ass:C11}
For any adversarial strategy $q \in \Aadv \cup
\Aadv^{\mathrm{env}}$ producing channel deviations bounded by
$\deltaadv$ (the Lemma~5b KL floor), the per-unit-time growth rate
of the proxy-truth gap is bounded:
\[
  \frac{d \epsgap}{dt} \;\leq\; \rho_g(\deltaadv) \;<\; \infty,
\]
where $\rho_g(\deltaadv)$ is a deployment-class constant depending
on $\mathcal{D}$'s operational parameters (channel-deviation
magnitudes, alignment-property smoothness in channel parameters)
but independent of $|P|$.
\end{assumption}

\textbf{Operational interpretation.}  The condition asserts that
during the detection window, the gap cannot grow arbitrarily fast.
Concretely: if channel deviations are bounded (else SPRT fires
faster), then the rate at which channel deviations translate into
proxy-truth gap is bounded.  This is structurally analogous to
Lemma~2's coordinate-Lipschitz condition (CL), but at the channel-
to-gap layer rather than the dimension-to-gap layer.

(C11) is new content; \citep[Microfoundation]{lasser2026micro}'s
``$T$ failure modes'' remark notes the bounded-Lipschitz-in-channel
property informally, but does not promote it to a theorem-level
condition.  Paper~10 surfaces it explicitly so the deployment
claim's Layer~2 intensivity is auditable.

\section{Statement of Lemma 6}

\begin{lemma}[$h_{\mathrm{detect}}$ intensivity]
\label{lem:6}
Under conditions (C5) continuous SPRT monitoring, (C7) bounded
co-evolution, and (C11) bounded per-unit-time gap-growth rate of
Theorem~\ref{thm:deployment_safety}, plus Lemma~\ref{lem:4}'s
SPRT lead-time tail bound:
\[
  h_{\mathrm{detect}}(\theta) \;\leq\; \rho_g(\deltaadv) \cdot
  T_{\mathrm{detect}}.
\]
With probability at least $1 - \beta'$ (Lemma~\ref{lem:5d}):
\[
  h_{\mathrm{detect}}(\theta) \;\leq\; \rho_g(\deltaadv) \cdot
  T_{\mathrm{cascade}}.
\]
$h_{\mathrm{detect}}(\theta)$ is intensive in $|P|$ over the
deployment class $\mathcal{D}$: $\rho_g(\deltaadv)$,
$T_{\mathrm{cascade}}$, and $\beta'$ are deployment-class constants.
\end{lemma}

\section{Proof}

\begin{proof}
We compose (C11) with Lemma~\ref{lem:4} to bound
$h_{\mathrm{detect}}$.

\emph{Step 1 (per-unit-time bound from C11).}  By
Assumption~\ref{ass:C11}, for any $q \in \Aadv \cup
\Aadv^{\mathrm{env}}$:
\[
  \frac{d \epsgap}{dt} \;\leq\; \rho_g(\deltaadv).
\]
$\rho_g(\deltaadv)$ is a deployment-class constant by (C11),
independent of $|P|$.

\emph{Step 2 (integrated bound over detection window).}  The gap
accumulated during the detection window:
\[
  h_{\mathrm{detect}}(\theta) \;=\;
  \int_{t_{\mathrm{violation}}}^{t_{\mathrm{violation}} +
  T_{\mathrm{detect}}} \frac{d \epsgap}{dt} \, dt
  \;\leq\; \rho_g(\deltaadv) \cdot T_{\mathrm{detect}}.
\]

\emph{Step 3 (apply Lemma 4 / Lemma 5d for tail bound).}
Lemma~\ref{lem:4} establishes the SPRT detection-time tail bound:
\[
  \Pr[T_{\mathrm{detect}} > T_{\mathrm{cascade}}] \;\leq\;
  \exp(-\kappa \cdot T_{\mathrm{cascade}} \cdot \deltaadv) \;=:\;
  \beta'.
\]
Therefore, with probability at least $1 - \beta'$,
$T_{\mathrm{detect}} \leq T_{\mathrm{cascade}}$, and:
\[
  h_{\mathrm{detect}}(\theta) \;\leq\; \rho_g(\deltaadv) \cdot
  T_{\mathrm{cascade}}.
\]

\emph{Step 4 (intensivity).}  Each constant is independent of
$|P|$ over $\mathcal{D}$:
\begin{itemize}
\item $\rho_g(\deltaadv)$: deployment-class constant by (C11).
\item $T_{\mathrm{cascade}} \geq \tau_{\mathrm{meta}}$:
  \citep[Phase~Redundancy]{lasser2026phase}'s metastable-lifetime
  lower bound (deployment-class constant per the (C7) bounded
  co-evolution caveat in Lemma~\ref{lem:4}'s appendix proof).
\item $\beta'$: deployment-class tail-bound constant from
  Lemma~\ref{lem:5d}'s composition.
\end{itemize}
Therefore $h_{\mathrm{detect}}(\theta) \leq \rho_g(\deltaadv) \cdot
T_{\mathrm{cascade}}$ is intensive in $|P|$ over $\mathcal{D}$.
$\Box$
\end{proof}

\section{Discussion}

\paragraph{Constant taxonomy.}
\begin{itemize}
\item $\rho_g(\deltaadv)$: \emph{deployment-policy-derived} (C11);
  measured by per-channel sensitivity-to-gap calibration.
\item $T_{\mathrm{cascade}}, \tau_{\mathrm{meta}}$: \emph{source-
  derived} from \citep[Phase~Redundancy]{lasser2026phase}'s
  metastable-lifetime scaling (with the (C7) bounded co-evolution
  qualification).
\item $\beta'$: \emph{composed constant} from Lemmas~5b, 4, 5d
  applied to the deployment's calibrated $\deltaadv$ and
  $\tau_{\mathrm{meta}}$.
\end{itemize}

\paragraph{Trade-off: tighter $\rho_g$ vs.\ wider regime.}
Deployments can tighten $\rho_g$ by:
\begin{itemize}
\item Restricting $g$'s class to alignment properties with low
  per-channel sensitivity (which also tightens (C6) Lipschitz).
\item Restricting $\Aadv$ to adversarial strategies whose per-step
  channel deviations have low gap-growth rate (effectively
  narrowing the detection class).
\end{itemize}
Each tightening produces a smaller $h_{\mathrm{detect}}$ but
potentially excludes some deployment configurations.  The
deployment-tooling specification calibrates $\rho_g$ in line with
the deployment's threat model.

\paragraph{Connection to (C11) versus per-step versus per-unit-time.}
(C11) is stated in continuous-time form $d\epsgap/dt$.  For
discrete-time deployments (per-step SPRT updates), the
corresponding form is $\rho_g^{\mathrm{step}}$, with
$h_{\mathrm{detect}} \leq \rho_g^{\mathrm{step}} \cdot
N_{\mathrm{detect}}$ where $N_{\mathrm{detect}}$ is the number of
SPRT steps to detection.  The continuous-time and discrete-time
forms are equivalent under standard step-time normalization.
Operators choose the formulation matching their deployment's SPRT
implementation.

\paragraph{Why (C11) is needed beyond (C6) Lipschitz alignment.}
(C6) bounds the Lipschitz constant of $g$ in the gap quantity
$\epsgap$.  (C11) bounds the rate at which $\epsgap$ itself can
grow under adversarial strategies in the detection class.  These
are different bounds: (C6) is alignment-property-side; (C11) is
gap-dynamics-side.  Both are needed for $h_{\mathrm{detect}}$
intensivity.

\section{Connection to Theorem 1 Layer 2}

Lemma~6 supplies the intensivity argument that
Theorem~\ref{thm:deployment_safety}'s Layer~2 statement asserts:
\[
  |g(\Ttruth) - g(\Pproxy)| \;\leq\; \mathrm{Lip}(g) \cdot
  h_{\mathrm{detect}}(\theta).
\]

After applying Lemma~6:
\[
  |g(\Ttruth) - g(\Pproxy)| \;\leq\; K_{\mathrm{Lip}} \cdot
  \rho_g(\deltaadv) \cdot T_{\mathrm{cascade}}.
\]
The right-hand side is the product of $|P|$-independent constants
over $\mathcal{D}$, with confidence at least $1 - \beta'$
(from Lemma~5d's tail composition).

This gives the explicit Layer~2 capability-magnitude-independent
bound:
\[
  |g(\Ttruth) - g(\Pproxy)| \;\leq\; C^{\mathrm{detect}}_*
  \quad \text{w.p. at least } 1 - \beta',
\]
where $C^{\mathrm{detect}}_* = K_{\mathrm{Lip}} \cdot
\rho_g(\deltaadv) \cdot T_{\mathrm{cascade}}$ is a deployment-class
constant.

\section{What this proof does and does not establish}

\textbf{Establishes:}
\begin{itemize}
\item A formal continuous-time per-unit-time gap-growth-rate bound
  (C11) as a deployment-class assumption.
\item Integrated $h_{\mathrm{detect}}$ bound combining (C11) with
  Lemma~4's tail bound and Lemma~5d's composition.
\item Explicit constant $C^{\mathrm{detect}}_*$ in deployment-class
  terms.
\item Connection to Theorem~1 Layer~2's intensivity claim.
\end{itemize}

\textbf{Does not establish (deferred):}
\begin{itemize}
\item Tightness of $\rho_g(\deltaadv)$.  Operators may calibrate
  tighter bounds for restricted threat models or tighter alignment
  properties.
\item Existence of deployment configurations satisfying (C11).
  Verified empirically by the deployment-tooling specification
  (a new audit, possibly Audit~7).
\item Numerical value of $\rho_g(\deltaadv)$.  Calibrated
  deployment-by-deployment.
\item Structural derivation of (C11) from finer source-paper
  machinery.  Currently asserted as a deployment-class condition,
  analogous to (C7) bounded co-evolution and SA1 HHI surrogate
  adequacy.
\end{itemize}

\section{Specific verification questions for codex review}

\begin{enumerate}
\item Is the continuous-time formulation of (C11) ($d\epsgap/dt$)
  the right formalization for Theorem~1's Layer~2?  Or should it
  be discrete-time (per-SPRT-step)?  The proof notes equivalence
  under step-time normalization, but Theorem~1's actual
  application context (continuous monitoring vs.\ discrete event
  ledger) might prefer one or the other.

\item Is $h_{\mathrm{detect}}$'s definition as the time-integrated
  gap growth correct?  Specifically, should it be:
  (a) $\int (d\epsgap/dt) dt$ (gap accumulated during detection
  window), or
  (b) $\sup_t \epsgap(t)$ (worst-case instantaneous gap)?
  The proof uses (a) but Theorem~1's application might prefer (b).

\item Is the (C11) condition correctly distinguished from (C6)
  bounded-Lipschitz alignment?  The discussion notes (C6) is
  alignment-property-side and (C11) is gap-dynamics-side.  Is
  this distinction watertight, or does (C11) reduce to a
  consequence of (C6) under some structural assumptions?

\item Is the connection to Lemma~5d's tail bound rigorous?
  Specifically, the argument is that with probability $1 - \beta'$,
  $T_{\mathrm{detect}} \leq T_{\mathrm{cascade}}$, so
  $h_{\mathrm{detect}} \leq \rho_g \cdot T_{\mathrm{cascade}}$.
  Is this the right way to translate the tail bound into the
  integrated gap bound, or does it lose information about the
  per-trajectory gap evolution?

\item Should (C11) be promoted to a theorem-level condition (C11)
  of Theorem~1, analogous to the (C6)-(C10) conditions added in
  earlier integrations?  Or can it remain a Lemma~6-internal
  assumption without theorem-level surfacing?
\end{enumerate}

\end{document}
